\documentclass[11pt]{article}
\usepackage{amsmath, amssymb, amsthm}
\usepackage[margin=1in]{geometry}
\newtheorem{proposition}{Proposition}
\newtheorem{lemma}{Lemma}
\title{\texttt{qcutelm\_vlt12}: efficient (pyramid) training for a\\
shared-weight recursive sandwich}
\author{qcute session notes}
\date{}

\begin{document}
\maketitle

\section{Setup and notation}

Let the model have $n$ levels, $i = 0, \dots, n-1$, with fixed local
compression factors $K_0, \dots, K_{n-1}$ and shared hidden width
$d$. Level $0$'s input is the raw byte sequence $b_{1:L_0}$
($L_0 = \texttt{context\_len}$); level $i>0$'s input is level
$i-1$'s code sequence $c^{(i-1)}_{1:L_i}$, where
\[
  L_i = L_0 \Big/ \prod_{j<i} K_j, \qquad
  n^{(i)}_{\text{blk}} = L_i / K_i \quad \text{(number of blocks at level $i$).}
\]

Three shared functions (Config.\texttt{untie\_levels=False}; all live in
$\mathbb{R}^d$, or $\mathbb{R}^{d_{\text{code}}}$ for codes, per
Config.\texttt{code\_dim}):
\begin{itemize}
  \item $E_\theta: \mathbb{R}^{T \times d} \to \mathbb{R}^{T \times d}$ — causal
        encoder (one shared weight set, applied at every level).
  \item $D_\theta: \mathbb{R}^{T \times d} \to \mathbb{R}^{T \times d}$ — causal
        decoder (one shared weight set).
  \item $F_\theta: \mathbb{R}^{T \times d_{\text{code}}} \to \mathbb{R}^{T \times
        d_{\text{code}}}$ — causal code-forecaster (\texttt{codelm}, one shared
        weight set).
\end{itemize}
and one shared readout $Q_\theta: \mathbb{R}^d \to \mathbb{R}^{d_{\text{code}}}$
(\texttt{code\_pre}) composed with a quantizer $\mathrm{quant}(\cdot)$
(BSQ/FSQ/IFSQ/identity — a fixed, parameter-free or STE-gradient function,
not part of the FLOP/pyramid argument below).

\section{Currently implemented: sequential per-level training}

\texttt{qcutelm\_vlt12.py}'s \texttt{forward()} (as of this note) computes,
for $i = 0, \dots, n-1$ in order:
\begin{align}
  h^{(i)}_a &= E_\theta\big(\mathrm{embed}_i(\mathrm{seq}_i)\big)
    && \text{[ENCODE, cost $\propto |{\rm seq}_i| = L_i$]} \label{eq:encode}\\
  c^{(i)} &= \mathrm{quant}\big(Q_\theta(h^{(i)}_{a,\,\text{block-ends}})\big)
    && \text{[readout, $n^{(i)}_{\text{blk}}$ positions]} \\
  \hat c^{(i)} &= F_\theta(c^{(i)})
    && \text{[FORECAST, cost $\propto n^{(i)}_{\text{blk}}$]} \label{eq:forecast}\\
  h^{(i)}_b &= D_\theta\big(\mathrm{subst}(\mathrm{embed}_i(\mathrm{seq}_i),\ \hat c^{(i)})\big)
    && \text{[DECODE, cost $\propto L_i$]} \label{eq:decode}\\
  \mathrm{seq}_{i+1} &:= c^{(i)}
\end{align}
i.e. \textbf{three separate self-attention passes per level} — one call to
$E_\theta$, one to $F_\theta$, one to $D_\theta$ — $3n$ total calls per
training step (here $n=3$: 9 separate stacks/step, matching this session's
own overhead-bound diagnosis for \texttt{qcutelm\_vlt11}, which has the
identical sequential structure just with untied weights). Total attention
FLOPs (dominant term, dense attention $O(T^2 d)$ per call, or
$O(T \cdot w \cdot d)$ windowed with window $w$):
\begin{equation}
  \mathrm{Cost}_{\text{sequential}} \;=\; \sum_{i=0}^{n-1}
    \Big[\underbrace{L_i}_{\text{encode}} + \underbrace{n^{(i)}_{\text{blk}}}_{\text{forecast}}
    + \underbrace{L_i}_{\text{decode}}\Big] \cdot O(w\, d)
  \;=\; O(w d) \sum_{i=0}^{n-1} \big(2 L_i + L_i/K_i\big).
  \label{eq:seq-cost}
\end{equation}
Weight-sharing changes \emph{which parameters} $E_\theta$/$D_\theta$/$F_\theta$
denote across the sum (same $\theta$ every term) but not the number of
\emph{calls} — this is a parameter-count and (via the gradient-isolation
argument in \texttt{qcutelm\_vlt12.py}'s docstring) a training-signal-pooling
win, not yet a wall-clock win. The wall-clock win requires \S\ref{sec:pyramid}.

\section{The pyramid reformulation}
\label{sec:pyramid}

\subsection{Key observation}

Because $E_\theta$ (resp.\ $D_\theta$, $F_\theta$) is the \emph{same
function} at every level, and every level's input sequence is causally
determined purely by the (already-computed) level below it, the encoder's
hidden state at pyramid position $(i, t)$ — level $i$, position $t$ — is a
\emph{fixed function of causal history up to that node}:
\begin{equation}
  h^{(i)}_a[t] \;=\; \big(E_\theta \text{ applied causally}\big)\big(c^{(i-1)}_{1:t}\big),
  \label{eq:node-fn}
\end{equation}
independent of anything about levels $>i$ (nothing above ever feeds
back down) and independent of \emph{which later reader} consults
$h^{(i)}_a[t]$ (there is exactly one reader per node in this
architecture — its own block's $Q_\theta$ readout — but the point
generalizes: node values never depend on the read side). This is
\texttt{docs/fifo\_v2.md} \S1's pyramid observation, transplanted from
FIFO's fixed byte-merge nodes to \texttt{vlt12}'s learned-code nodes —
it applies here specifically \emph{because} $E_\theta$ is shared (in
\texttt{qcutelm\_vlt11}'s untied version, ``the function applied at level
$i$'' and ``the function applied at level $i{+}1$'' are literally
different functions with different weights, so there is no single
pyramid node value to share between them).

\subsection{One-pass construction}

Define the pyramid as the disjoint union of every level's node set,
$\mathcal{P} = \bigsqcup_{i=0}^{n-1} \{(i, t) : 1 \le t \le L_i\}$, total
size
\begin{equation}
  |\mathcal{P}| \;=\; \sum_{i=0}^{n-1} L_i \;=\; L_0 \Big(1 + \tfrac{1}{K_0}
    + \tfrac{1}{K_0 K_1} + \cdots\Big) \;<\; L_0 \cdot \tfrac{K}{K-1}
  \label{eq:pyramid-size}
\end{equation}
for compression factors all $\ge K$ — a convergent geometric series, so
$|\mathcal{P}| = O(L_0)$, not $O(n L_0)$: the pyramid is at most a small
constant factor larger than the base sequence (e.g.\ $K{=}4$ gives
$|\mathcal{P}| < \tfrac{4}{3} L_0$), regardless of $n$.

Embed every node into $\mathbb{R}^d$ (bytes via the $8$-bit projection,
codes via $\mathrm{embed}_i$ — all landing in the same shared space per
Config.\texttt{code\_dim}), concatenate into one sequence of length
$|\mathcal{P}|$, and run \emph{one} call to $E_\theta$ with an attention
mask $M$ such that node $(i,t)$ may attend to node $(i',t')$ iff
\begin{equation}
  M\big[(i,t), (i',t')\big] = 1 \iff
  \text{the byte-span of } (i',t') \text{ lies entirely within}
  \text{the causal history of } (i,t).
  \label{eq:mask}
\end{equation}
Concretely: $(i,t)$ may attend to every $(i,t')$ with $t' \le t$ (its own
level's causal history — recovers \eqref{eq:encode}'s within-level
causality) and every $(i-1, t')$ with $t' \le t \cdot K_{i-1}$ (the finer
level's history up to the same real-time point — recovers the fact that
$c^{(i-1)}$'s block ending at $t$ is itself a function of $b$-history up
to $t \cdot K_{i-1} \cdots K_0$). This is a \emph{block-lower-triangular},
not fully dense, mask — its cost is discussed in \S\ref{sec:cost}.

\begin{lemma}
Given mask \eqref{eq:mask}, node $(i,t)$'s output under the one-pass
pyramid attention equals $h^{(i)}_a[t]$ from the sequential
computation \eqref{eq:encode}.
\end{lemma}
\begin{proof}[Proof sketch]
By induction on $i$. Base case $i=0$: mask restricts $(0,t)$ to attend
only to $(0, t' \le t)$, which is exactly causal self-attention over
raw bytes — identical to \eqref{eq:encode} at $i=0$. Inductive step:
assume nodes at level $i-1$ already equal their sequential values (by
hypothesis). Node $(i,t)$ attends to level-$(i-1)$ nodes up to the
corresponding causal cutoff and level-$i$ nodes up to $t$; since
$Q_\theta \circ \mathrm{quant}$ applied to level-$(i-1)$'s (already
correct) hidden states reproduces $c^{(i-1)}$ exactly as in the
sequential path, and $E_\theta$ is the same function/weights either way,
$(i,t)$'s output matches \eqref{eq:encode} at level $i$.
\end{proof}

An identical construction applies to $D_\theta$ (decode nodes, reading
the same finer-level history plus the substituted forecast at block
starts) and $F_\theta$ (forecast nodes, reading only same-level code
history) — three pyramid passes total instead of $3n$ sequential ones,
each pass handling every level's nodes for that role at once.

\subsection{Cost comparison}
\label{sec:cost}

Dense attention over $|\mathcal{P}|$ nodes with the mask
\eqref{eq:mask} costs $O(|\mathcal{P}|^2 d)$ in the worst case (a node
can attend to $O(|\mathcal{P}|)$ others), matching
\texttt{docs/fifo\_v2.md} \S1's own cost note
($\mathrm{pyramid} \sim (2 W_{\max})^2$ for FIFO's analogous
construction). With \emph{windowed} attention (this file's existing
$\texttt{attn\_window}$ convention, extended to the pyramid mask: each
node attends only to the most recent $w$ nodes at its own level, plus
the most recent $w \cdot K_{i-1}$ real-time-equivalent span one level
down), cost is $O(|\mathcal{P}| \cdot w \cdot d) = O(L_0 \cdot w \cdot d)$
by \eqref{eq:pyramid-size} — \emph{independent of $n$}, versus the
sequential cost \eqref{eq:seq-cost}'s $\sum_i (2L_i + L_i/K_i)$, which is
also $O(L_0 \cdot w \cdot d)$ up to the same geometric-series constant.

\textbf{The pyramid's real saving is not asymptotic FLOPs} (both are
$O(L_0)$ up to constants) \textbf{— it is call count.} Sequential:
$3n$ separate kernel-launched self-attention calls per step (this
session's own overhead-bound diagnosis: \texttt{qcutelm\_vlt11} at
$n=3$ measured 4--9$\times$ fewer FLOPs/step than \texttt{bytelm} yet ran
\emph{slower}, because MPS pays real per-call dispatch overhead
independent of each call's FLOP count). Pyramid: \textbf{3 calls total},
regardless of $n$. On a throughput regime where per-call overhead
dominates FLOPs (exactly the regime measured this session), this is the
lever that matters — collapsing $3n \to 3$ calls is a genuine,
measurable wall-clock win even though the total attention FLOPs are
asymptotically unchanged.

\subsection{Reading off level-specific quantities}

Once the pyramid pass is computed, level $i$'s forecast
$\hat c^{(i)}$, code-match loss, and NTP loss are all pure \emph{index
operations} into $\mathcal{P}$'s already-computed hidden states —
gather the $n^{(i)}_{\text{blk}}$ nodes belonging to level $i$, apply
the (also shared) head, done. No further attention. This mirrors
\texttt{docs/fifo\_v2.md} \S1's own claim (``cheap — indexing + small
per-node MTP heads, no further attention'') verbatim, now applied to
\texttt{vlt12}'s learned-code levels instead of FIFO's sampled
compositions.

\section{Status}

\textbf{Not implemented} in \texttt{qcutelm\_vlt12.py} as of this note —
the file's \texttt{forward()} runs the sequential form
(\S2), verified correct and matching v11's mechanics exactly (including
the detach-teacher-force gradient-isolation property, verified to survive
weight-sharing — see the module docstring). The pyramid form (\S3) is the
documented target for a follow-up implementation: building the block
mask \eqref{eq:mask} for windowed multi-level attention, and the gather
step in \S3.4, are the two pieces of real new code required. This
mirrors \texttt{docs/fifo\_v2.md}'s own convention of separating ``design
sketch, with a correctness argument'' from ``implemented'' — kept
explicit here rather than conflated.

\end{document}
